\chapter{قابلیت‌ها، محدودیت‌ها و ادامهٔ کار}

نسخهٔ اول گفتگو، حافظهٔ صریح، مدیریت کار، سرچ اطلاعات تازه، خواندن 
\lr{HTML}
 عمومی، بررسی کد پروژه، تغییر فایل کوچک با تأیید، تاریخچهٔ موقت و دستور مستقیم محلی دارد. جدول 
\lr{docs/CAPABILITIES.md}
 فرق پیاده‌سازی و شاهد تست را مشخص می‌کند.

اجرا نشده‌اند: فرمان دلخواه سیستم، تست خودکار کد، ورودی تصویر و صوت و 
\lr{PDF}
، کنترل مرورگر و ورود، کار زمان‌بندی‌شده، یادآوری، سرویس چندکاربره، رابط گرافیکی، جستجوی برداری و 
\lr{cloud}
 نامحدود. این‌ها انتخاب آینده هستند و قابلیت پنهان فعلی نیستند. کاربر برای این نسخه روش ترمینالی را انتخاب کرده است.

فایل 
\lr{PROJECT\_STATUS.md}
 معماری، تصمیم محدوده، نتیجهٔ تست، محدودیت حریم خصوصی، ساختار گزارش و کار اختیاری آینده را ثبت می‌کند. همکار بعدی اول آن را بخواند و بعد چهار راهنمای اولیه را بررسی کند. گزارش دوزبانه و قابل فهم و با نام مالک پروژه باشد. پس از تغییر معنادار، نتیجهٔ تست و آموزش را به‌روز کنید؛ قابلیت را فقط از متن پرامپت نتیجه نگیرید.

قدم بعدی مناسب پس از استفادهٔ واقعی، جمع کردن چند نمونهٔ خطای واقعی با حذف راز و رفع کوچک‌ترین مشکل مشخص است. پچ فایل بزرگ یا رابط گرافیکی می‌تواند بعداً طراحی شود. قبل از آزمودن جریان فعلی در کار روزمره، چند چارچوب و لایهٔ اتوماسیون اضافه نکنید.

نام دقیق نویسنده و شش نشانی ارتباطی در تنظیم متمرکز گزارش قرار دارند. فونت دانلودی، تصویر بررسی، 
\lr{PDF}
 و دادهٔ شخصی محلی و خارج از 
\lr{Git}
 می‌مانند. منبع عمومی و آموزش و گزارش پیش از 
\lr{commit}
 دستی قابل بررسی‌اند. در این کار 
\lr{commit}
 یا 
\lr{push}
 یا انتشار یا استقرار انجام نشده است.
